\documentclass[a4paper,11pt]{ltjsarticle}
\usepackage{geometry}
\geometry{top=25mm,bottom=25mm,left=25mm,right=25mm}
\usepackage{listings}
\usepackage{booktabs}
\usepackage{fancyhdr}
\usepackage{hyperref}
\usepackage{xcolor}
\usepackage{amsmath}
\usepackage{longtable}
\usepackage{array}

\lstset{
  basicstyle=\ttfamily\small,
  keywordstyle=\color{blue}\bfseries,
  commentstyle=\color{green!50!black},
  stringstyle=\color{red},
  numbers=left, numberstyle=\tiny\color{gray}, numbersep=5pt,
  frame=single, breaklines=true, tabsize=4, showstringspaces=false, captionpos=b
}

\pagestyle{fancy}
\fancyhf{}
\fancyhead[L]{計算機科学実験及演習3 ハードウェア — ユーザーズマニュアル}
\fancyfoot[C]{\thepage}
\renewcommand{\headrulewidth}{0.4pt}

\begin{document}

\begin{titlepage}
\centering
\vspace*{3cm}
{\LARGE 計算機科学実験及演習3 ハードウェア}\\[1cm]
{\Huge\bfseries ユーザーズマニュアル}\\[0.5cm]
{\LARGE SIMPLE/B 拡張プロセッサシステム}\\[3cm]
\begin{tabular}{rl}
学籍番号 & （学籍番号を記入） \\[0.5em]
氏名     & （氏名を記入） \\[0.5em]
提出日   & 2026年6月5日 \\
\end{tabular}
\vfill
京都大学 工学部情報学科 計算機科学コース
\end{titlepage}

\tableofcontents
\clearpage


% =====================================================================
\section{システム概要}

本プロセッサシステムは SIMPLE アーキテクチャ（ver 4.0）に基づく16ビットプロセッサであり、
以下の拡張を施している。

\begin{itemize}
\item 4サイクル実行方式（基本5サイクルから25\%高速化）
\item ISA拡張: BAL（サブルーチン呼出）、BR（間接ジャンプ）、ADDI（即値加算）
\item ハードウェアタイマ（プリスケーラ付き16ビットカウンタ）
\item I/Oアドレス空間（最大16ポート、6ポート実装済み）
\item 8桁7セグメントLED表示（4モード切替）
\item LED出力（8ビット）
\end{itemize}

\subsection{対象ハードウェア}

\begin{table}[h]
\centering
\caption{ハードウェア仕様}
\begin{tabular}{ll}
\toprule
項目 & 内容 \\
\midrule
FPGAボード & PowerMedusa MU500-RX/RK \\
FPGA & Altera Cyclone IV EP4CE30F23I7N \\
クロック & 20 MHz 発振器 \\
主記憶 & Block RAM 4096語 $\times$ 16ビット \\
\bottomrule
\end{tabular}
\end{table}


% =====================================================================
\section{操作方法}

\subsection{起動と実行}

\begin{enumerate}
\item FPGAボードの電源を投入する
\item Quartus PrimeからFPGAにプログラムを書き込む（\texttt{.sof}ファイル）
\item リセットボタン（\texttt{rst\_n}、負論理）を押して初期化する
\item 実行ボタン（\texttt{exec}、正論理）を押すとプログラムが実行開始する
\item HLT命令でプロセッサが停止する。再実行するにはリセット後に再度execを押す
\end{enumerate}

\subsection{入力デバイス}

\begin{table}[h]
\centering
\caption{入力デバイス一覧}
\begin{tabular}{lll}
\toprule
デバイス & 接続 & 備考 \\
\midrule
DIPスイッチ A (SW26) & \texttt{sw[7:0]} & 負論理 (ON=0)、I/Oポート0で読出し \\
DIPスイッチ B (SW27) & \texttt{sw[15:8]} & 負論理 (ON=0)、I/Oポート1で読出し \\
プッシュスイッチ     & \texttt{pushsw[3:0]} & 負論理 (押下=0)、I/Oポート4で読出し \\
ロータリースイッチ   & \texttt{rotsw[3:0]} & I/Oポート5で読出し、表示モード切替兼用 \\
\bottomrule
\end{tabular}
\end{table}

\subsection{出力デバイス}

\begin{table}[h]
\centering
\caption{出力デバイス一覧}
\begin{tabular}{lll}
\toprule
デバイス & 接続 & 備考 \\
\midrule
7セグメントLED & \texttt{seg[7:0]}, \texttt{sel[7:0]} & 8桁HEX表示、ダイナミック点灯 \\
LED            & \texttt{led[7:0]} & I/Oポート4で書込み \\
\bottomrule
\end{tabular}
\end{table}

\subsection{7セグメント表示モード}

ロータリースイッチの下位2ビットで表示内容を切り替える。

\begin{table}[h]
\centering
\caption{7セグ表示モード}
\begin{tabular}{ccl}
\toprule
ロータリー値 & モード & 表示内容 \\
\midrule
0 & OUT値表示 & OUTポート0/1の値を8桁HEX表示（デフォルト） \\
1 & PC表示   & プログラムカウンタの値（デバッグ用） \\
2 & レジスタ表示 & r0\textasciitilde r3の下位8ビットを各2桁で表示 \\
3 & タイマ表示 & タイマCOUNT値 \\
\bottomrule
\end{tabular}
\end{table}


% =====================================================================
\section{命令セットリファレンス}

\subsection{命令形式}

全命令は16ビット固定長であり、4種類の形式を持つ。

\begin{table}[h]
\centering
\caption{命令形式}
\begin{tabular}{llp{7cm}}
\toprule
形式 & ビット配置 & 用途 \\
\midrule
(a) & \texttt{11|Rs(3)|Rd(3)|op3(4)|d(4)} & 演算, シフト, IN/OUT, HLT \\
(b) & \texttt{op1(2)|Ra(3)|Rb(3)|d(8)} & LD (op1=00), ST (op1=01) \\
(c) & \texttt{10|op2(3)|Rb(3)|d(8)} & LI, BAL, BR, ADDI, B \\
(d) & \texttt{10|111|cond(3)|d(8)} & Bcc (条件分岐) \\
\bottomrule
\end{tabular}
\end{table}

\subsection{レジスタ}

8本の汎用レジスタ \texttt{r0}\textasciitilde\texttt{r7}（各16ビット）を持つ。
慣例として \texttt{r7} をBAL命令のリンクレジスタとして使用する。

\subsection{演算命令}

\begin{longtable}{llp{6cm}l}
\caption{演算命令一覧} \label{tab:alu_ops} \\
\toprule
ニーモニック & op3 & 動作 & フラグ \\
\midrule
\endfirsthead
\toprule
ニーモニック & op3 & 動作 & フラグ \\
\midrule
\endhead
\texttt{ADD Rd, Rs} & 0000 & \texttt{r[Rd] = r[Rd] + r[Rs]} & S,Z,C,V \\
\texttt{SUB Rd, Rs} & 0001 & \texttt{r[Rd] = r[Rd] - r[Rs]} & S,Z,C,V \\
\texttt{AND Rd, Rs} & 0010 & \texttt{r[Rd] = r[Rd] \& r[Rs]} & S,Z \\
\texttt{OR Rd, Rs}  & 0011 & \texttt{r[Rd] = r[Rd] | r[Rs]} & S,Z \\
\texttt{XOR Rd, Rs} & 0100 & \texttt{r[Rd] = r[Rd] \^{} r[Rs]} & S,Z \\
\texttt{CMP Rd, Rs} & 0101 & \texttt{r[Rd] - r[Rs]}（フラグのみ） & S,Z,C,V \\
\texttt{MOV Rd, Rs} & 0110 & \texttt{r[Rd] = r[Rs]} & S,Z \\
\midrule
\texttt{SLL Rd, d}  & 1000 & 左論理シフト & S,Z,C \\
\texttt{SLR Rd, d}  & 1001 & 左循環シフト & S,Z \\
\texttt{SRL Rd, d}  & 1010 & 右論理シフト & S,Z,C \\
\texttt{SRA Rd, d}  & 1011 & 右算術シフト & S,Z,C \\
\midrule
\texttt{IN Rd, port} & 1100 & \texttt{r[Rd] = input[port]} & --- \\
\texttt{OUT Rs, port} & 1101 & \texttt{output[port] = r[Rs]} & --- \\
\texttt{HLT}         & 1111 & 実行停止 & --- \\
\bottomrule
\end{longtable}

\subsection{メモリアクセス命令}

\begin{table}[h]
\centering
\caption{ロード/ストア命令}
\begin{tabular}{lp{8cm}}
\toprule
ニーモニック & 動作 \\
\midrule
\texttt{LD Ra, d(Rb)} & \texttt{r[Ra] = mem[r[Rb] + sign\_ext(d)]} \\
\texttt{ST Ra, d(Rb)} & \texttt{mem[r[Rb] + sign\_ext(d)] = r[Ra]} \\
\bottomrule
\end{tabular}
\end{table}

変位 \texttt{d} は8ビット符号付き（$-128$\textasciitilde$+127$）。

\subsection{即値・分岐命令}

\begin{table}[h]
\centering
\caption{即値・分岐命令}
\begin{tabular}{llp{7cm}}
\toprule
ニーモニック & op2 & 動作 \\
\midrule
\texttt{LI Rb, d}     & 000 & \texttt{r[Rb] = sign\_ext(d)} \\
\texttt{B d}           & 100 & \texttt{PC = PC + 1 + sign\_ext(d)} \\
\texttt{BE/BLT/BLE/BNE d} & 111 & 条件成立時 \texttt{PC = PC + 1 + sign\_ext(d)} \\
\bottomrule
\end{tabular}
\end{table}

\subsubsection{条件分岐の条件}

\begin{table}[h]
\centering
\caption{条件分岐コード}
\begin{tabular}{lll}
\toprule
ニーモニック & cond & 条件 \\
\midrule
\texttt{BE}  & 000 & Z = 1 （等しい） \\
\texttt{BLT} & 001 & S $\oplus$ V = 1 （符号付き小なり） \\
\texttt{BLE} & 010 & Z $\vee$ (S $\oplus$ V) = 1 （符号付き以下） \\
\texttt{BNE} & 011 & Z = 0 （等しくない） \\
\bottomrule
\end{tabular}
\end{table}

\subsection{拡張命令}

\begin{table}[h]
\centering
\caption{ISA拡張命令}
\begin{tabular}{llp{7cm}l}
\toprule
ニーモニック & op2 & 動作 & フラグ \\
\midrule
\texttt{BAL Rb, d} & 001 & \texttt{r[Rb] = PC; PC = PC + 1 + sign\_ext(d)} & --- \\
\texttt{BR Rb}     & 010 & \texttt{PC = r[Rb]} & --- \\
\texttt{ADDI Rb, d}& 011 & \texttt{r[Rb] = r[Rb] + sign\_ext(d)} & S,Z,C,V \\
\bottomrule
\end{tabular}
\end{table}

\subsubsection{BAL/BR の使い方}

\texttt{BAL} はサブルーチン呼び出しに、\texttt{BR} はサブルーチンからの復帰に使用する。

\begin{lstlisting}[language={},caption={サブルーチン呼出の例}]
; メインプログラム
        BAL  r7, sub       ; sub を呼び出し (r7 = 戻りアドレス)
        OUT  r0             ; sub から戻った後に実行される
        HLT

; サブルーチン
sub:    LI   r0, 42
        BR   r7             ; r7 の値 (=呼出元の次の命令) に戻る
\end{lstlisting}


% =====================================================================
\section{I/Oポートマップ}

IN/OUT 命令の第2オペランドでポート番号を指定する。
省略時はポート0が使用される（後方互換）。

\begin{table}[h]
\centering
\caption{I/Oポートマップ}
\begin{tabular}{clll}
\toprule
ポート & IN (読み出し) & OUT (書き込み) & アセンブラ記法 \\
\midrule
0 & DIPスイッチ A (8bit) & 7セグ出力 下位16bit & \texttt{IN r0, 0} \\
1 & DIPスイッチ B (8bit) & 7セグ出力 上位16bit & \texttt{IN r0, 1} \\
2 & タイマ COUNT値 (16bit)  & タイマ CTRL レジスタ & \texttt{OUT r0, 2} \\
3 & タイマ overflow (bit0)  & タイマ LIMIT 値      & \texttt{OUT r0, 3} \\
4 & プッシュスイッチ (4bit) & LED出力 (8bit)       & \texttt{OUT r0, 4} \\
5 & ロータリースイッチ (4bit) & (予約)             & \texttt{IN r0, 5} \\
\bottomrule
\end{tabular}
\end{table}


% =====================================================================
\section{タイマの使い方}

\subsection{レジスタ}

\begin{table}[h]
\centering
\caption{タイマレジスタ}
\begin{tabular}{lclp{6cm}}
\toprule
レジスタ & ポート & R/W & 説明 \\
\midrule
COUNT    & IN 2  & R & 現在のカウント値（16ビット） \\
CTRL     & OUT 2 & W & bit0: 動作(1)/停止(0), bit1: リセット, bit[15:8]: プリスケーラ \\
LIMIT    & OUT 3 & W & カウント上限値 \\
overflow & IN 3  & R & bit0: 上限到達フラグ (読出しで自動クリア) \\
\bottomrule
\end{tabular}
\end{table}

\subsection{プリスケーラ}

カウント周波数 $= 20\,\mathrm{MHz} \div (\mathrm{PRESCALER} + 1)$。

\begin{itemize}
\item PRESCALER = 0: 20 MHz（50 ns 間隔）
\item PRESCALER = 19: 1 MHz（1 $\mu$s 間隔）
\item PRESCALER = 199: 100 kHz（10 $\mu$s 間隔）
\end{itemize}

\subsection{使用例: 時間計測}

\begin{lstlisting}[language={},caption={タイマによる実行時間計測}]
; タイマ初期化
        LI   r1, 0xFF       ; LIMIT = 0x00FF (上位=0)
        OUT  r1, 3           ; LIMIT レジスタに設定

; カウント開始 (プリスケーラ=0, 20MHz)
        LI   r1, 1           ; CTRL: bit0=1 (動作開始)
        OUT  r1, 2

; ... (計測対象の処理) ...

; カウント値読出し
        IN   r0, 2           ; r0 = COUNT値
        OUT  r0              ; 7セグに表示

; タイマ停止
        LI   r1, 0
        OUT  r1, 2           ; CTRL = 0 (停止)
\end{lstlisting}

\subsection{使用例: 一定時間待ち}

\begin{lstlisting}[language={},caption={タイマによる一定時間待ち}]
; LIMIT=9999, プリスケーラ=199 で 1秒タイマ
; (100kHz / 10000 = 10Hz → 100ms, LIMIT=9999でCOUNT 0~9999)
        LI   r1, 0x0F        ; LIMIT下位 (仮)
        ; (...上位ビットの設定はLI+SLL+ORで構築...)
        OUT  r1, 3            ; LIMIT設定

; CTRL = プリスケーラ199(=0xC7), EN=1 → 0xC701
        ; (...16ビット定数の構築...)
        OUT  r1, 2            ; タイマ動作開始

; overflow待ちループ
wait:   IN   r0, 3            ; overflow読出し
        LI   r1, 1
        AND  r0, r1           ; bit0を取得
        BE   wait             ; overflow=0 なら待ち継続
        ; overflow=1 で脱出 (読出し時に自動クリア済み)
\end{lstlisting}


% =====================================================================
\section{アセンブラの使い方}

\subsection{概要}

\texttt{simple\_asm.py} は全命令に対応したPythonアセンブラである。
ラベル、コメント、PC相対分岐のアドレス自動計算をサポートする。

\subsection{実行方法}

\begin{lstlisting}[language={},caption={アセンブラの実行}]
# 基本的な使い方
python3 simple_asm.py input.asm -o output.hex

# リスティング付きで実行
python3 simple_asm.py input.asm -o output.hex -l

# 出力ファイル名を省略すると input.hex が生成される
python3 simple_asm.py input.asm
\end{lstlisting}

\subsection{アセンブリ言語の文法}

\begin{lstlisting}[language={},caption={アセンブリ言語の文法}]
; 行コメント (セミコロン以降が無視される)
label:  MNEMONIC  operand1, operand2   ; インラインコメント

; ラベル定義: 行頭の識別子 + コロン
; 数値リテラル: 10 (10進), 0x0A (16進), 0b1010 (2進)
; 負の即値: -1, -128
; レジスタ: r0 ~ r7 (大文字/小文字可)
\end{lstlisting}

\subsection{命令記法一覧}

\begin{table}[h]
\centering
\caption{アセンブラ記法}
\small
\begin{tabular}{ll}
\toprule
記法 & 例 \\
\midrule
\texttt{ADD Rd, Rs}      & \texttt{ADD r0, r1} \\
\texttt{LD Ra, d(Rb)}    & \texttt{LD r0, 5(r1)} / \texttt{LD r0, -3(r2)} \\
\texttt{ST Ra, d(Rb)}    & \texttt{ST r0, 0(r1)} \\
\texttt{LI Rb, imm}      & \texttt{LI r0, 0x80} / \texttt{LI r1, -1} \\
\texttt{B label}          & \texttt{B loop} \\
\texttt{Bcc label}        & \texttt{BLT inner} / \texttt{BNE done} \\
\texttt{BAL Rb, label}    & \texttt{BAL r7, func} \\
\texttt{BR Rb}            & \texttt{BR r7} \\
\texttt{ADDI Rb, imm}     & \texttt{ADDI r0, 1} / \texttt{ADDI r3, -1} \\
\texttt{IN Rd [, port]}   & \texttt{IN r0} / \texttt{IN r0, 2} \\
\texttt{OUT Rs [, port]}  & \texttt{OUT r0} / \texttt{OUT r1, 4} \\
\texttt{HLT}              & \texttt{HLT} \\
\bottomrule
\end{tabular}
\end{table}


% =====================================================================
\section{プログラムの開発手順}

\subsection{シミュレーション}

\begin{lstlisting}[language={},caption={iverilog によるシミュレーション}]
cd simple/

# アセンブル
python3 simple_asm.py myprog.asm -o myprog.hex -l

# コンパイル (INIT_FILE を指定)
iverilog -DINIT_FILE=\"myprog.hex\" -o sim.out \
    simple_tb.v simple_cpu.v alu.v shifter.v regfile.v ram_sim.v

# 実行
vvp sim.out
\end{lstlisting}

\subsection{FPGA実装}

\begin{enumerate}
\item アセンブラでhexファイルを生成
\item hexファイルからMIFファイルを作成し、\texttt{run.mif}として保存
\item Quartus Primeでプロジェクトを開き、全Verilogモジュールを追加:
      \texttt{simple\_top.v}, \texttt{simple\_cpu.v}, \texttt{alu.v},
      \texttt{shifter.v}, \texttt{regfile.v}, \texttt{ram.v},
      \texttt{seg7dec.v}, \texttt{timer.v}
\item ピンアサインを設定（\texttt{.qsf}ファイル）
\item SDCファイルでクロック制約を設定
\item コンパイル → Programmer でFPGAに書き込み
\end{enumerate}


% =====================================================================
\section{サンプルプログラム}

\subsection{Hello World（即値加算と表示）}

\begin{lstlisting}[language={},caption={最初のプログラム}]
; hello.asm - 3 + 5 = 8 を7セグに表示
        LI   r0, 3          ; r0 = 3
        LI   r1, 5          ; r1 = 5
        ADD  r0, r1         ; r0 = r0 + r1 = 8
        OUT  r0              ; 7セグに "0008" 表示
        HLT                  ; 停止
\end{lstlisting}

\subsection{DIPスイッチ読取り}

\begin{lstlisting}[language={},caption={DIPスイッチの値を表示}]
; dipsw.asm - DIPスイッチの値を読み取って表示
loop:   IN   r0, 0          ; DIPスイッチAを読取り
        IN   r1, 1          ; DIPスイッチBを読取り
        SLL  r1, 8          ; 上位8ビットにシフト
        OR   r0, r1         ; r0 = B:A (16ビット)
        OUT  r0              ; 7セグに表示
        B    loop            ; 無限ループ
\end{lstlisting}

\subsection{LED制御}

\begin{lstlisting}[language={},caption={LEDをカウントアップ}]
; ledcount.asm - LEDを0~255でカウントアップ
        LI   r0, 0          ; カウンタ
loop:   OUT  r0, 4          ; LEDに出力 (ポート4)
        OUT  r0              ; 7セグにも表示
        ; 遅延 (ソフトウェアループ)
        LI   r1, 0
delay:  ADDI r1, 1
        BNE  delay           ; r1 がオーバーフロー(0)するまで
        ADDI r0, 1           ; カウンタ++
        B    loop
\end{lstlisting}

\subsection{バブルソート}

\texttt{sort.asm} を参照。16要素の配列をバブルソートし、結果を7セグに表示する
デモプログラム。


% =====================================================================
\section{トラブルシューティング}

\begin{table}[h]
\centering
\caption{よくある問題と対処法}
\small
\begin{tabular}{p{5cm}p{8cm}}
\toprule
症状 & 対処法 \\
\midrule
execを押しても動かない & リセットボタンを押してから再度execを押す \\
7セグが何も表示しない & ロータリースイッチがモード0になっているか確認。OUT命令がプログラムに含まれているか確認 \\
シミュレーションで期待と違う出力 & リスティング (\texttt{-l}) でアドレスを確認。分岐先ラベルの綴りを確認 \\
アセンブルエラー「即値が範囲外」 & 8ビット即値は $-128$\textasciitilde$+127$（符号付き）または $0$\textasciitilde$255$（符号なし）。16ビット定数はLI+SLL+ORで構築する \\
DIPスイッチの値が反転する & 負論理。ハードウェア側で反転済み（\texttt{$\sim$sw}）なので、ソフトウェアではそのまま使用可能 \\
\bottomrule
\end{tabular}
\end{table}


\end{document}
